\documentclass[10pt,a4paper]{amsart}
\usepackage{amsmath,amssymb,amsthm,graphicx}
% Retained preamble for the Stacks Project chapter snapshot in this directory.
% The chapter's own preamble.tex is a copy of the shared Stacks Project preamble; the block
% below is sliced MECHANICALLY, line for line, from it. Nothing was added, removed or
% redefined; the wrapper class, the xr-hyper/hyperref wiring and the externaldocument list
% are omitted because the assembler supplies the document class.
%
% lines taken from preamble.tex: 3, 5, 9, 11, 12, 156, 157, 158, ... (51 lines total)
%
% package coverage restored from the chapter's own preamble (2026-10-13): xy, xy-2cell

\usepackage[all]{xy}
\xyoption{2cell}
\UseAllTwocells

\newenvironment{reference}{\comment}{\endcomment}
\newenvironment{slogan}{\comment}{\endcomment}
\newenvironment{history}{\comment}{\endcomment}
\theoremstyle{plain}
\newtheorem{theorem}[subsection]{Theorem}
\newtheorem{proposition}[subsection]{Proposition}
\newtheorem{lemma}[subsection]{Lemma}
\theoremstyle{definition}
\newtheorem{definition}[subsection]{Definition}
\newtheorem{example}[subsection]{Example}
\newtheorem{exercise}[subsection]{Exercise}
\newtheorem{situation}[subsection]{Situation}
\theoremstyle{remark}
\newtheorem{remark}[subsection]{Remark}
\newtheorem{remarks}[subsection]{Remarks}
\numberwithin{equation}{subsection}
\def\lim{\mathop{\mathrm{lim}}\nolimits}
\def\colim{\mathop{\mathrm{colim}}\nolimits}
\def\Spec{\mathop{\mathrm{Spec}}}
\def\Hom{\mathop{\mathrm{Hom}}\nolimits}
\def\Ext{\mathop{\mathrm{Ext}}\nolimits}
\def\SheafHom{\mathop{\mathcal{H}\!\mathit{om}}\nolimits}
\def\SheafExt{\mathop{\mathcal{E}\!\mathit{xt}}\nolimits}
\def\Sch{\mathit{Sch}}
\def\Mor{\mathop{\mathrm{Mor}}\nolimits}
\def\Ob{\mathop{\mathrm{Ob}}\nolimits}
\def\Sh{\mathop{\mathit{Sh}}\nolimits}
\def\NL{\mathop{N\!L}\nolimits}
\def\CH{\mathop{\mathrm{CH}}\nolimits}
\def\proetale{{pro\text{-}\acute{e}tale}}
\def\etale{{\acute{e}tale}}
\def\QCoh{\mathit{QCoh}}
\def\Ker{\mathop{\mathrm{Ker}}}
\def\Im{\mathop{\mathrm{Im}}}
\def\Coker{\mathop{\mathrm{Coker}}}
\def\Coim{\mathop{\mathrm{Coim}}}
\DeclareMathSymbol{\boxtimes}{\mathbin}{AMSa}{"02}
\def\QCohstack{\mathcal{QC}\!\mathit{oh}}
\def\Cohstack{\mathcal{C}\!\mathit{oh}}
\def\Spacesstack{\mathcal{S}\!\mathit{paces}}
\def\Quotfunctor{\mathrm{Quot}}
\def\Hilbfunctor{\mathrm{Hilb}}
\def\Curvesstack{\mathcal{C}\!\mathit{urves}}
\def\Polarizedstack{\mathcal{P}\!\mathit{olarized}}
\def\Complexesstack{\mathcal{C}\!\mathit{omplexes}}
\def\Pic{\mathop{\mathrm{Pic}}\nolimits}
\def\Picardstack{\mathcal{P}\!\mathit{ic}}
\def\Picardfunctor{\mathrm{Pic}}
\def\Deformationcategory{\mathcal{D}\!\mathit{ef}}
\usepackage{cleveref}
\crefname{originallemma}{original lemma}{original lemmas}
\Crefname{originallemma}{Original lemma}{Original lemmas}
\makeatletter
\newlabel{definition-picard-functor}{{0B9L}{source}{}{sourcepointer-1}{}}
\newlabel{definition-picard-functor@cref}{{[originallemma][0][]0B9L}{[1][0][]source}{}{}{}}
\makeatother
\crefname{originallemma}{original lemma}{original lemmas}
\Crefname{originallemma}{Original lemma}{Original lemmas}
\makeatletter
\newlabel{descent-proposition-fpqc-descent-quasi-coherent}{{023T}{source}{}{sourcepointer-2}{}}
\newlabel{descent-proposition-fpqc-descent-quasi-coherent@cref}{{[originallemma][0][]023T}{[1][0][]source}{}{}{}}
\makeatother
\crefname{originallemma}{original lemma}{original lemmas}
\Crefname{originallemma}{Original lemma}{Original lemmas}
\makeatletter
\newlabel{lemma-flat-geometrically-connected-fibres}{{0B9M}{source}{}{sourcepointer-3}{}}
\newlabel{lemma-flat-geometrically-connected-fibres@cref}{{[originallemma][0][]0B9M}{[1][0][]source}{}{}{}}
\makeatother
\title{Split exact sequence for the Picard functor with a section}
\author{The Stacks Project Authors}
\date{stacks-project:pic.tex:local-snapshot-20261012; GNU FDL 1.2 or later}
\begin{document}
\maketitle
\noindent\textit{Editorial note.} Statement span is lines 1027-1038 of sources/stacks-pic/pinned-original-source.tex; proof span is lines 1040-1093, the complete author proof (52 body lines).\par\medskip
\noindent\textit{Editorial note.} Navigation stubs: 3 cited-but-not-included labels are supplied as non-load-bearing pointers with their pinned tags where the snapshot tag index has them.\par\medskip
\noindent\textit{Editorial note.} extra\_packages is empty on purpose: the retained preamble already loads xy with options and a second declaration clashes in the assembler wrapper.\par\medskip
\noindent\textit{Editorial note.} Text is retained verbatim, including the author's own typos and placeholders.\par\medskip
\noindent\textit{Editorial note.} Queue index 6 of queue\_pic.json (pb=2253, stmt 1027-1038, proof 1040-1093).\par\medskip
\noindent\textit{Editorial note.} stacks-pic wave-1 (pi01): fresh adjudication accept (2026-10-14); pool per standing queue hypercovering -> pic (SMC reassigned to Zhen per Chief 08:5x); pic.tex 27 labels/23 with proof judged 21 accept / 2 reject / 4 pb=0 env-dup rows excluded; spans cite-free and reference/slogan/history-free so no patchenv needed. Named omission: Author notes '(small detail omitted)' in the descent application - named tiny check; the rescaling and f\_ijk = 1 arguments are in-span.\par\medskip
\noindent\textit{Allowed standard prerequisites (editorial).}\begin{itemize}
\item Explicit prerequisite boundary. The cited inputs are prerequisites; the argument in the retained span is not replaced by any cited result.
\end{itemize}\medskip
\ifdefined\hypertarget\hypertarget{sourcepointer-1}{}\fi
\noindent\textit{Source reference.} Original lemma 0B9L: sources/stacks-pic/pinned-original-source.tex, source lines 957--968. Cited by the retained statement or proof but not part of the collected spans; supplied only as a navigation pointer. The argument does not depend on its content.\par\medskip
\ifdefined\hypertarget\hypertarget{sourcepointer-2}{}\fi
\noindent\textit{Source reference.} Original lemma 023T: sources/stacks-pic/pinned-original-source.tex, source lines 1093--1093. Cited by the retained statement or proof but not part of the collected spans; supplied only as a navigation pointer. The argument does not depend on its content.\par\medskip
\ifdefined\hypertarget\hypertarget{sourcepointer-3}{}\fi
\noindent\textit{Source reference.} Original lemma 0B9M: sources/stacks-pic/pinned-original-source.tex, source lines 979--987. Cited by the retained statement or proof but not part of the collected spans; supplied only as a navigation pointer. The argument does not depend on its content.\par\medskip
\begin{lemma}
\label{lemma-flat-geometrically-connected-fibres-with-section}
Let $f : X \to S$ be as in Definition \ref{definition-picard-functor}.
Assume $f$ has a section $\sigma$ and that
$\mathcal{O}_T \to f_{T, *}\mathcal{O}_{X_T}$ is an isomorphism
for all $T \in \Ob((\Sch/S)_{fppf})$. Then
$$
0 \to \Pic(T) \to \Pic(X_T) \to \Picardfunctor_{X/S}(T) \to 0
$$
is a split exact sequence with splitting given by
$\sigma_T^* : \Pic(X_T) \to \Pic(T)$.
\end{lemma}

\begin{proof}
Denote $K(T) = \Ker(\sigma_T^* : \Pic(X_T) \to \Pic(T))$.
Since $\sigma$ is a section of $f$ we see that $\Pic(X_T)$ is the direct
sum of $\Pic(T)$ and $K(T)$.
Thus by Lemma \ref{lemma-flat-geometrically-connected-fibres} we see that
$K(T) \subset \Picardfunctor_{X/S}(T)$ for all $T$. Moreover, it is clear
from the construction that $\Picardfunctor_{X/S}$ is the sheafification
of the presheaf $K$. To finish the proof it suffices to show that
$K$ satisfies the sheaf condition for fppf coverings which we do
in the next paragraph.

\medskip\noindent
Let $\{T_i \to T\}$ be an fppf covering. Let $\mathcal{L}_i$ be
elements of $K(T_i)$ which map to the same elements of $K(T_i \times_T T_j)$
for all $i$ and $j$. Choose an isomorphism
$\alpha_i : \mathcal{O}_{T_i} \to \sigma_{T_i}^*\mathcal{L}_i$
for all $i$. Choose an isomorphism
$$
\varphi_{ij} :
\mathcal{L}_i|_{X_{T_i \times_T T_j}}
\longrightarrow
\mathcal{L}_j|_{X_{T_i \times_T T_j}}
$$
If the map
$$
\alpha_j|_{T_i \times_T T_j} \circ
\sigma_{T_i \times_T T_j}^*\varphi_{ij} \circ
\alpha_i|_{T_i \times_T T_j} :
\mathcal{O}_{T_i \times_T T_j} \to \mathcal{O}_{T_i \times_T T_j}
$$
is not equal to multiplication by $1$ but some $u_{ij}$, then we can scale
$\varphi_{ij}$ by $u_{ij}^{-1}$ to correct this. Having done this, consider
the self map
$$
\varphi_{ki}|_{X_{T_i \times_T T_j \times_T T_k}} \circ
\varphi_{jk}|_{X_{T_i \times_T T_j \times_T T_k}} \circ
\varphi_{ij}|_{X_{T_i \times_T T_j \times_T T_k}}
\quad\text{on}\quad
\mathcal{L}_i|_{X_{T_i \times_T T_j \times_T T_k}}
$$
which is given by multiplication by some regular function $f_{ijk}$
on the scheme $X_{T_i \times_T T_j \times_T T_k}$.
By our choice of $\varphi_{ij}$ we see that the pullback of
this map by $\sigma$ is equal to multiplication by $1$. By
our assumption on functions on $X$, we see that $f_{ijk} = 1$.
Thus we obtain a descent datum for the fppf covering
$\{X_{T_i} \to X\}$. By
Descent, Proposition \ref{descent-proposition-fpqc-descent-quasi-coherent}
there is an invertible $\mathcal{O}_{X_T}$-module $\mathcal{L}$
and an isomorphism $\alpha : \mathcal{O}_T \to \sigma_T^*\mathcal{L}$
whose pullback to $X_{T_i}$ recovers $(\mathcal{L}_i, \alpha_i)$
(small detail omitted). Thus $\mathcal{L}$ defines an object
of $K(T)$ as desired.
\end{proof}

\end{document}
